\documentclass[a4paper]{article}

\usepackage{graphicx}
\usepackage{amsmath,amssymb}
\usepackage{minted}
\usepackage{multirow}
\usepackage{float}
\usepackage{todonotes}
\usepackage{forest}
\usepackage{xstring}
\usepackage{xcolor}
\usepackage[most]{tcolorbox}
\usepackage{xparse}
\usepackage{natbib}

\usetikzlibrary{graphs,graphdrawing}
\usegdlibrary{layered}

% for external graphviz
\input{/home/donkarlo/Dropbox/repo/nd_language_ai_project/src/nd_language_ai/formal/computer/programming/tex/latex/code_clips/external_graphviz.tex}

% for tags
\input{/home/donkarlo/Dropbox/repo/nd_language_ai_project/src/nd_language_ai/formal/computer/programming/tex/latex/part/control_sequence/kind/semtag/semtag.tex}

% for links
\input{/home/donkarlo/Dropbox/repo/nd_language_ai_project/src/nd_language_ai/formal/computer/programming/tex/latex/code_clips/seehere.tex}

\title{Bayesian theories og brain}
\date{}

\begin{document}
   \maketitle
   Bayesian approaches to brain function investigate the capacity of the nervous system to operate in situations of uncertainty in a fashion that is close to the optimal prescribed by Bayesian statistics. \url{https://en.wikipedia.org/wiki/Bayesian_approaches_to_brain_function}
   \section{Evidences}
      \subsection{Particle filter}
         Kutschireiter et al. propose a Neural Particle Filter (NPF), a sampling-based nonlinear Bayesian filter intended as a biologically plausible model of perceptual inference \cite{kutschireiter-2017-nonlinear-bayesian-filtering-and-learning-a-neuronal-dynamics-for-perception}.
         In the model, particles representing samples from the posterior distribution are directly associated with the activity of neurons in a recurrent neural network.
         The resulting neuronal dynamics can integrate noisy sensory information over time and across multiple sensory modalities while performing approximate Bayesian state estimation.


      \subsection{Kalman filter}
         Human sensorimotor control operates under substantial sensory and motor uncertainty, and Bayesian decision theory provides a framework for explaining how the central nervous system performs estimation and control under such uncertainty \cite{wolpert-2007-probabilistic-models-in-human-sensorimotor-control}. Experimental evidence suggests that humans combine multiple sensory sources approximately according to maximum-likelihood estimation and incorporate prior knowledge to form maximum a posteriori estimates.
         The review also discusses evidence that the nervous system uses Kalman-filter-like processes to estimate dynamically changing states of the body and environment.

         %\bibliographystyle{plainnat}
         %\bibliography{/home/donkarlo/Dropbox/repo/nd_shared_research/bibliography/bibliography.bib}
\end{document}